\documentclass{beamer}
\usetheme{Montpellier}

\usepackage[utf8]{inputenc}
\usepackage[T1]{fontenc}

\usepackage{hyperref}

\usepackage[
backend=bibtex,
style=authortitle-comp,
]{biblatex}
\addbibresource{presentation.bib}
\DeclareFieldFormat{url}{\newline\scriptsize{\url{#1}}}

\usepackage{enumitem}
\setlist[description]{style=nextline}

\AtBeginSection[]{
  \begin{frame}
    \vfill
    \centering
    \begin{beamercolorbox}[sep=8pt,center,shadow=true,rounded=true]{title}
      \usebeamerfont{title}\insertsectionhead\par%
    \end{beamercolorbox}
    \vfill
  \end{frame}
}

\usepackage{graphicx}
\graphicspath{{./resources/}}

\usepackage{minted}
\definecolor{gql}{RGB}{229,53,171}
\setminted[gql]{
  autogobble=true,
  fontfamily=tt,
  fontsize=\footnotesize,
  frame=leftline,
  framerule=1pt,
  rulecolor=gql,
}
\definecolor{ts}{RGB}{49,120,198}
\setminted[ts]{
  autogobble=true,
  fontfamily=tt,
  fontsize=\footnotesize,
  frame=leftline,
  framerule=1pt,
  rulecolor=ts,
}
\setminted[json]{
  autogobble=true,
  fontfamily=tt,
  fontsize=\footnotesize,
  frame=leftline,
  framerule=1pt,
  rulecolor=black,
}

\title{GraphQL:\@Thinking in Resolvers}
\author{Sébastien Lavoie-Courchesne}
\date{February 23rd 2022}

\begin{document}
\section*{Introduction}
\addtocounter{section}{1}
\begin{frame}
  \titlepage{}
\end{frame}

\begin{frame}
  \begin{center}
    Don't know how your clients are using your APIs?

    \vspace{2em}
    Just return everything\ldots{}the right way!
  \end{center}
\end{frame}


\begin{frame}
  \frametitle{Sébastien Lavoie-Courchesne}
  \begin{itemize}
  \item Architect for the Catalog Group at AppDirect
  \item Member of the GraphQL Working Group
  \item Developed most of the initial infrastructure we use
  \end{itemize}
\end{frame}


\begin{frame}
  \frametitle{Contents}
  \tableofcontents
\end{frame}


\begin{frame}
  \frametitle{Technologies used in the examples}
  \begin{minipage}{0.4\linewidth}
    \begin{flushright}
      GraphQL
      
    \vspace{2em}
    NodeJS
    
    \vspace{2em}
    TypeScript
    
    \vspace{2em}
    ApolloServer
    \end{flushright}
  \end{minipage}
  \begin{minipage}{0.45\linewidth}
    \begin{center}
      \includegraphics[height=3em]{GraphQL}

      \includegraphics[height=3em]{NodeJS}
    
      \includegraphics[height=3em]{TypeScript}
    
      \includegraphics[height=3em]{ApolloServer}
    \end{center}
  \end{minipage}
\end{frame}
\note{
  Colourized frame on code examples

  pink for GraphQL
  
  blue for TS
  
  black for JSON
}


\begin{frame}
  \frametitle{Chirper}
  
  A chirping service

  \vspace{1em}
  Create a user

  \vspace{1em}
  Send chirps

  \vspace{1em}
  Reply to chirps
\end{frame}


\section{GraphQL}
\begin{frame}
  \frametitle{What is GraphQL?}  From the GraphQL
  Specification\cite{graphql-spec}:
  \begin{quote}
    GraphQL is a query language designed to build client applications
    by providing an intuitive and flexible syntax and system for
    describing their data requirements and interactions.
  \end{quote}
\end{frame}


\begin{frame}
  \frametitle{But What is GraphQL?}

  \begin{itemize}
  \item Specification
  \item Query language
  \item Communication over the network
  \item Server and client requirements
  \end{itemize}
\end{frame}


\begin{frame}
  \frametitle{Why GraphQL?}
  Only request and receive the fields you want:
  \begin{itemize}
  \item Smaller payloads
  \item Descriptive language
  \item Documented
  \end{itemize}
\end{frame}
\note{
  Often compared to REST and gRPC
  
  Mostly used over HTTP, but not always

  Particularly useful when a server doesn't know about all its clients
}


\begin{frame}
  \frametitle{How does it work?}
  Server defines the schema it serves using the GraphQL language

  \vspace{1em}
  Clients can read the schema and send queries to the server using the GraphQL language, optionally passing in variables
\end{frame}
\note{
  typically one query per request, multiple queries in a single request are supported
}


\begin{frame}
  \frametitle{Operations}
  \begin{description}
  \item[Query] a read-only fetch, analogous to REST GET
  \item[Mutation] a write followed by a fetch, analogous to REST
    POST/PATCH/PUT/DELETE
  \item[Subscription] a long-lived request that fetches data in
    response to source events
  \end{description}
\end{frame}
\note{
  Subscription is less often used and less documented, multiple ways to serve it (WebSocket, Kafka, RabbitMQ, etc)
}


\begin{frame}
  \frametitle{HTTP Communication}
  Most common use of GraphQL
  
  HTTP GET or POST to a specific URL
  
  Traditional HTTP headers/cookies

  Returns status code 200

  \vspace{1em}
  GET \texttt{query} parameter contains the query as a string
  
  POST JSON Body includes:

  \hspace{1em}\texttt{query}, \texttt{variables}, \texttt{operationName}
\end{frame}
\note{
  query as a string

  variables as a JSON object
  
  operationName is used to differentiate multiple requests sent in a single call on the client side

  errors are returned within the response body, clients must parse it
}


\begin{frame}
  \frametitle{Server}
  Servers handle requests by:

  \begin{itemize}
  \item parsing query string
  \item building execution plan
  \item executing plan if valid
  \item returning result
  \end{itemize}
\end{frame}


\begin{frame}[fragile]
  \frametitle{Example Request}
  \begin{minted}{json}
    {
      "query": "
        query ($id: ID!) {
          chirp(id: $id) {
            contents
          }
        }
      ",
      "variables": {
        "id": "1234"
      }
    }
  \end{minted}
\end{frame}


\begin{frame}[fragile]
  \frametitle{Example Response}
  \begin{minted}{json}
    {
      "data": {
        "chirp": {
          "contents": "a chirp"
        }
      },
      "errors": null
    }
  \end{minted}
\end{frame}


\section{GraphQL Schema Language}
\begin{frame}[fragile=singleslide]
  \frametitle{Primitives}
  \mintinline{gql}{scalar} represent a single
  value. Predefined: ID, String, Int, Float, Boolean.
  \begin{minted}{gql}
    """
    Date and Time represented using the ISO 8601 standard
    """
    scalar DateTime

    """
    Email represented as a String
    Format validated when used as input
    """
    scalar Email
  \end{minted}
\end{frame}

\begin{frame}[fragile]
  \frametitle{Primitives}
  \mintinline{gql}{enum} represents an enumeration of possible values.

  Serialized as strings.

  \begin{minted}{gql}
    """
    Cardinal directions
    """
    enum Direction {
      NORTH
      EAST
      SOUTH
      WEST
    }
  \end{minted}
\end{frame}


\begin{frame}[fragile]
  \frametitle{Modifiers}
  \mintinline{gql}{!} to denote a non-nullable field
  \begin{minted}{gql}
    """
    Non-nullable identifier
    """
    id: ID!
  \end{minted}
\end{frame}
\note{
  nullability can be used to indicate a missing value, returning null for a non-nullable type generates an error
}

\begin{frame}[fragile]
  \frametitle{Modifiers}
  \mintinline{gql}{[]} to denote a list of elements, around the type
  \begin{minted}{gql}
    """A nullable list of nullable strings"""
    a: [String]
    
    """A non-nullable list of nullable strings"""
    b: [String]!
    
    """A nullable list of non-nullable strings"""
    c: [String!]
    
    """A non-nullable list of non-nullable strings"""
    d: [String!]!
  \end{minted}
\end{frame}
\note{
  an empty list is non-null

  good practice to have non-nullable elements within the list
}


\begin{frame}[fragile]
  \frametitle{Objects}
  \mintinline{gql}{input} used to represent a list of named and typed input fields, can only contain primitives and other input types
  \begin{minted}{gql}
    """Input used to create a new chirp"""
    input ChirpInput {
      """Contents of the chirp"""
      contents: String!
    }
  \end{minted}
\end{frame}
\note{
}


\begin{frame}[fragile]
  \frametitle{Objects}
  \mintinline{gql}{type} used to represent a list of named
  and typed output fields
  \begin{minted}{gql}
    """A Chirp"""
    type Chirp {
      """Identifier of the chirp"""
      id: ID!

      """Contents of the chirp"""
      contents(format: ContentsFormat = PLAIN): String!
    }
  \end{minted}
\end{frame}
\note{
  arguments can be added to fields to allow specifying more
  information about what should be returned, format, locale, page for
  collections, etc

  only input objects and primitives can be used as arguments, input
  fields cannot have arguments
}


\begin{frame}[fragile]
  \frametitle{Abstractions}
  \mintinline{gql}{interface} used to represent a list of named and typed output fields, implemented by types
  \begin{minted}{gql}
    """An identified and queryable node within the graph"""
    interface Node {
      """Unique non-nullable identifier of the node"""
      id: ID!
    }

    """An identifiable Chirp"""
    type Chirp implements Node {
      """Identifier of the chirp"""
      id: ID!

      """Contents of the chirp"""
      contents: String!
    }
  \end{minted}
\end{frame}
\note{
  any number of types can implement the interface
}


\begin{frame}[fragile]
  \frametitle{Abstractions}
  \mintinline{gql}{union} to represent a collection of possible output types
  \begin{minted}{gql}
    """A union of possible errors that can happen when creating a new chirp"""
    union ChirpUsageError = EmptyContents | TooLongContents

    """Contents were empty"""
    type EmptyContentsError implements UsageError {
      message: String!
    }

    """Contents were too long"""
    type TooLongContents implements UsageError {
      message: String!
      length: Int!
      maxLength: Int!
    }
  \end{minted}
\end{frame}
\note{
  unions can be combined with interfaces to create more powerful abstractions

  done by adding the same interface to every type in the union
}


\begin{frame}[fragile]
  \frametitle{Directives}
  Annotations on the graph, prefixed with \mintinline{gql}{@}, can have arguments.

  Predefined:
  \begin{description}
  \item [\mintinline{gql}{@deprecated(reason: String!)}]
    indicate to the client the element is deprecated, along with a reason
  \item [\mintinline{gql}{@skip(if: Boolean!)}]
    don't return the field's value when the condition is false
  \item [\mintinline{gql}{@include(if: Boolean!)}]
    return the field's value only when the condition is true
  \item [\mintinline{gql}{@specifiedBy(url: String!)}]
    add formatting information for scalars, usually pointed to a standard or specification
  \end{description}
\end{frame}


\begin{frame}[fragile]
  \frametitle{Directives}
  Custom directives can also be created
  \begin{minted}{gql}
    # defining a custom directive to validate authorization
    # for output fields
    directive @authorization(scopes: [String!]!) on FIELD_DEFINITION

    type Mutation {
      chirp(input: ChirpInput!): ChirpPayload!
        @authorization(scopes: ["chirp.write"])
    }
  \end{minted}
\end{frame}


\section{GraphQL Query Language}
\begin{frame}[fragile]
  \frametitle{Description}
  Clients use the query language to request information from a GraphQL server.

  \vspace{1em}
  \begin{minted}{gql}
    query {
      chirp(id: "1234") {
        contents
        author { username }
      }
    }
  \end{minted}
\end{frame}


\begin{frame}[fragile]
  \frametitle{Variables}
  Variables can be passed to parameterize requests

  \vspace{1em}
  \begin{minted}{gql}
    query ($id: ID!) {
      chirp(id: $id) {
        contents
        author { username }
      }
    }
  \end{minted}

  \vspace{1em}
  \begin{minted}{json}
    {
      "query": "query($id){chirp(id:$id){contents author{username}}}",
      "variables": {
        "id": "1234"
      }
    }
  \end{minted}
\end{frame}
\note{
  client-side advantage: static string construction (faster runtime)

  server-side advantage: can cache query execution plan (faster responses)
}


\section{Resolver}
\begin{frame}[fragile]
  \frametitle{Description}
  A resolver is a function that is used to return the value for a specific field within a type.

  \vspace{1em}
  \begin{minted}{gql}
    type User { name: String! }
  \end{minted}

  \vspace{1em}
  \begin{minted}{ts}
    resolvers = {
      User: {
        name: () => 'Sebastien Lavoie-Courchesne'
      }
    }
  \end{minted}
\end{frame}


\begin{frame}
  \frametitle{Anatomy} Resolvers take in 4 parameters:
  \begin{description}
  \item[parent] The parent object
  \item[arguments] Any arguments passed to the field
  \item[context] The context generated by the server
  \item[info] Additional information about the request
  \end{description}
\end{frame}


\begin{frame}[fragile]
  \frametitle{Default Resolver}
  If no resolver is defined for a field:
  \begin{minted}{ts}
    function [Type.fieldName](
      parent: Type
    ): any {
      return parent[fieldName];
    }
  \end{minted}
\end{frame}
\note{
  if field is a value, return it

  if field is a function, execute it and return its output
}


\begin{frame}
  \frametitle{Parent}
  Whatever was returned by the resolver of the parent field, including properties that are not mapped to fields in the GraphQL schema

  \vspace{1em}
  \mintinline{ts}{null} if the parent type is Query, Mutation, or Subscription
\end{frame}


\begin{frame}[fragile]
  \frametitle{Arguments}
  Any arguments passed to the field.

  \begin{minted}{gql}
    type Query {
      chirps(first: Int, after: String, last: Int, before: String):
        ChirpConnection!
    }
  \end{minted}

  \begin{minted}{ts}
    interface ChirpsArguments {
      first: number;
      after: string;
      last: number;
      before: string;
    }

    function chirps(parent: null, arguments: ChirpsArguments) // ...
  \end{minted}
\end{frame}


\begin{frame}[fragile]
  \frametitle{Context}
  A context object generated by the ApolloServer

  \vspace{1em}
  \begin{minipage}{0.6\linewidth}
    \begin{minted}{ts}
      new ApolloServer({
        context({ req, res }) {
          return {
            user: req.user,
          };
        },
        dataSources() {
          return {
            users: new UserDataSource(),
            chirps: new ChirpDataSource(),
          };
        }
      });
    \end{minted}
  \end{minipage}
  \begin{minipage}{0.05\linewidth}
    $\Rightarrow{}$
  \end{minipage}
  \begin{minipage}{0.25\linewidth}
    \begin{minted}{ts}
      {
        user,
        dataSources: {
          users,
          chirps,
        }
      }
    \end{minted}
  \end{minipage}
\end{frame}
\note{
  executed on each request
  
  contains the result of the context function as well as a dataSources property containing the result of the dataSources function
}


\begin{frame}[fragile]
  \frametitle{Data Sources}
  \begin{minted}{ts}
    abstract class DataSource<TContext> {
      initialize?(config: DataSourceConfig<TContext>):
        void | Promise<void>;
    }
  \end{minted}

  \vspace{1em}
  Initialize is called on each incoming request

  \vspace{1em}
  Instantiate a new data source on each incoming request
\end{frame}
\note{
  a second incoming request will re-initialize the data source
}


\begin{frame}
  \frametitle{What to include?}
  Context:
  \begin{itemize}
  \item authentication information
  \item logger
  \item tracer + trace information
  \end{itemize}

  \vspace{1em}
  DataSources:
  \begin{itemize}
  \item repositories
  \item caches
  \item services
  \item clients to other services
  \end{itemize}
\end{frame}


\begin{frame}[fragile]
  \frametitle{info}
  \begin{minted}{ts}
    interface GraphQLResolveInfo {
      fieldName: string;
      fieldNodes: FieldNode[];
      returnType: GraphQLOutputType;
      parentType: GraphQLObjectType;
      path: Path;
      schema: GraphQLSchema;
      fragments: Record<string, FragmentDefinitionNode>;
      rootValue: any;
      operation: OperationDefinitionNode;
      variableValues: Record<string, any>;
    }
  \end{minted}
\end{frame}
\note{
  mostly extracted to other fields

  path can be used to return the path when returning errors to identify where the error occurred
}


\begin{frame}
  \frametitle{Scalars}
  Scalars can have additional parsing in the form of a 3 function object:

  \begin{description}
  \item[parseLiteral] Used to parse values passed in through the query's text itself
  \item[parseValue] Used to parse values passed through the variables
  \item[serialize] Used to serialize the domain value to a String
  \end{description}

  \vspace{1em}
  These can be used to add validation (correct format, specific length, etc) as well as map it to another type (Date, URL, etc).
\end{frame}

\begin{frame}
  \frametitle{Directives} Directives can have additional logic tied to them in the form of a Visitor pattern.

  \vspace{1em}
  Directives are not visible to the clients introspecting the schema.

  \vspace{1em}
  These can be used to add validation (correct format, specific length, etc) or transformations (lower/uppercase, rounding, etc)
\end{frame}

\section{Federation}
\begin{frame}
  \frametitle{Specification}
  Apollo Federation is an additional specification to specify subgraphs, gateways and schema registries.

  \vspace{1em}
  The gateway aggregates the subgraphs to expose a single GraphQL endpoint. Subgraphs can reference each other.
\end{frame}

\begin{frame}[fragile]
  \frametitle{Subgraph}
  Individual service exposing a GraphQL endpoint. Services must satisfy an additional schema containing:

  \begin{minted}{gql}
    scalar _Any
    scalar _FieldSet
    
    # a union of all types that use the @key directive
    union _Entity
    
    type _Service {
      sdl: String
    }
    
    extend type Query {
      _entities(representations: [_Any!]!): [_Entity]!
      _service: _Service!
    }
  \end{minted}
\end{frame}
\note{
  this is entirely handled by most GraphQL implementations
}


\begin{frame}[fragile]
  \frametitle{Subgraph directives}
  The following directives are available to allow referencing other types within other subgraphs:

  \begin{minted}{gql}
    directive @external on FIELD_DEFINITION
    directive @requires(fields: _FieldSet!) on FIELD_DEFINITION
    directive @provides(fields: _FieldSet!) on FIELD_DEFINITION
    directive @key(fields: _FieldSet!) repeatable on OBJECT | INTERFACE
    
    # not all implementations allow this one
    # there's an "extend" keyword as well
    directive @extends on OBJECT | INTERFACE
  \end{minted}
\end{frame}


\begin{frame}[fragile]
  \frametitle{Subgraph References}
  Subgraphs can reference types from another subgraph by extending them:

  \begin{minted}{gql}
    extend type User @key(fields: "id") {
      id: ID! @external

      chirps: [Chirp!]!
    }

    type Chirp @key(fields: "id") {
      id: ID!

      author: User!
    }
  \end{minted}
\end{frame}
\note{
  chirp service adds chirps field on the user type defined in the user service

  chirp service also returns a reference to the user that authored the chirp

  user service has no knowledge of chirps
}


\begin{frame}[fragile]
  \frametitle{Returning references}
  Return the type and values for all the fields in the key:
  \begin{minted}{json}
    {
      "__typename": "User",
      "id": "1234"
    }
  \end{minted}
\end{frame}
\note{
  typename is only required when returning a union/interface

  but a good practice to keep it
}


\begin{frame}[fragile]
  \frametitle{Resolving references}
  Special resolver to resolve an entity based on its key:
  \begin{minted}{ts}
    function __resolveReference(
      parent: { id: string },
      { dataSources: { users } }: Context
    ): Promise<User> {
      return users.getById(id);
    }
  \end{minted}
\end{frame}
\note{
  parent, context, info

  no arguments parameter
}


\begin{frame}[fragile]
  \frametitle{Testing in isolation}
  \texttt{\_\_resolveReference} can be tested by using the \texttt{\_entities} query:
  \begin{minted}{gql}
    query {
      _entities(representations: [
        {
          __typename: "User",
          id: "1234"
        }
      ]) {
        ... on User {
          username
        }
      }
    }
  \end{minted}
\end{frame}
\note{
  this is equivalent to the query sent by the gateway when it needs to resolve an entity by its key
}


\begin{frame}
  \frametitle{Gateway}

  Main responsabilities on server start and/or periodically:
  \begin{itemize}
  \item Validate schemas from each subgraph
  \item Aggregate schemas from each subgraph
  \item Expose single GraphQL endpoint
  \end{itemize}

  \vspace{1em}
  Main responsabilities on each request:
  \begin{itemize}
  \item Parse incoming requests
  \item Build query execution plan
  \item Execute plan, sending sub-requests to each subgraph
  \item Return aggregated result
  \end{itemize}
\end{frame}


\begin{frame}
  \frametitle{Schema Registry}
  A Schema Registry can be used to validate and aggregate schemas.

  \vspace{1em}
  The gateway then pulls the schema from the schema registry.
\end{frame}
\note{
  Apollo offers a paid online schema registry

  there's also an open source schema registry that you can deploy yourself

  Reduces load and potential failures on the gateway
}


\section{Best Practices}
\begin{frame}
  \frametitle{Federation}
  When working with federation across multiple teams:
  \begin{itemize}
  \item Define guidelines \& standards early
  \item Have consistent naming strategies
  \item Review schemas to ensure they follow guidelines \& standards
  \end{itemize}
\end{frame}


\begin{frame}[fragile]
  \frametitle{Node interface}
  Every identifiable type should implement this node interface
  \begin{minted}{gql}
    """An identifiable node within the graph"""
    interface Node {
      """Unique identifier of the node"""
      id: ID!
    }
  \end{minted}

  And the server can expose a node query to fetch any node within the graph
  \begin{minted}{gql}
    type Query {
      """
      A global node query to get any node within the
      graph by its unique identifier
      """
      node(id: ID!): Node
    }
  \end{minted}
\end{frame}


\begin{frame}[fragile]
  \frametitle{Connection Pattern}
  For paginated queries, use the connection pattern\cite{Connection-Spec}

  \begin{minted}{gql}
    type PageInfo {
      hasPreviousPage: Boolean!
      hasNextPage: Boolean!
      startCursor: String
      endCursor: String
    }

    type ChirpConnection {
      pageInfo: PageInfo!
      edges: [Chirp!]!
    }

    type Query {
      chirps(first: Int, after: String, last: Int, before: String):
        ChirpConnection!
    }
  \end{minted}
\end{frame}
\note{
  Originally designed by Relay

  many client libraries understand this pattern and can automate the pagination
}


\begin{frame}[fragile]
  \frametitle{Mutation Pattern}
  Have a consistent pattern for mutations

  \begin{minted}{gql}
  type Mutation {
    chirp(input: ChirpInput!): ChirpPayload!
  }

  input ChirpInput {
    contents: String!
  }

  type ChirpPayload {
    chirp: Chirp
    errors: [ChirpUsageError!]
  }

  union ChirpUsageError = EmptyContents | ContentsTooLong
  \end{minted}
\end{frame}
\note{
  adding/removing fields is easier with an input object than multiple arguments

  returning a payload with either the created/updated entity or a list of errors

  also easier to update in the future
}


\begin{frame}
  \frametitle{Persisted Queries}
  Servers can allow creation of persisted queries.

  \vspace{1em}
  Client sends static query text to server (separate endpoint)\\
  Server parses, generates query plan, saves everything, return id

  \vspace{1em}
  For each query, client sends id of the persisted query and variables
\end{frame}
\note{
  smaller payloads on each request because you don't send the full query text

  faster execution on the server
}


\begin{frame}
  \frametitle{Rate Limiting}
  Rate limit by allocating points to each client

  \vspace{1em}
  Calculate the cost of each request in points:
  \begin{itemize}
  \item querying a single node: 1 point
  \item querying a list: 1 point + 1 point per element
  \item executing a mutation: x points
  \end{itemize}
\end{frame}
\note{
  multiple different algorithms to calculate the cost

  exact points are up to you

  can also specify points independently through directives
}


\section{Conclusion}
\begin{frame}
  \frametitle{Conclusion}
\end{frame}

\nocite{*}
\begin{frame}
  \frametitle{Specifications}
  \printbibliography[keyword=spec]
\end{frame}

\begin{frame}
  \frametitle{Material}
  \printbibliography[keyword=material]
\end{frame}

\end{document}
